\begin{table}[!h]
\centering
\caption{ERPT por signo del shock, para cada shock estructural ($h=12$)}
\centering
\fontsize{9}{11}\selectfont
\begin{threeparttable}
\begin{tabular}[t]{lcccc}
\toprule
Shock & ERPT$^{+}$ & $F^{+}$ & ERPT$^{-}$ & $F^{-}$\\
\midrule
Shock MP Fed & 0.202 & 0.1 & 0.181 & 0.0\\
Shock petróleo & -0.599 & 0.0 & -0.072 & 1.4\\
EBP & 0.145 & 1.9 & -0.438 & 0.3\\
Términos interc. & 0.164 & 0.2 & -0.044 & 2.9\\
EMBI & 0.643 & 0.3 & 0.017 & 5.0\\
\addlinespace
Cambiario (TCN) & 0.336 & 0.5 & -0.009 & 3.9\\
\bottomrule
\end{tabular}
\begin{tablenotes}
\item \textit{Note: } 
\item Cada shock del punto 6 se parte en pieza positiva y negativa; el ERPT de cada signo se estima por IV con esa pieza como instrumento y la otra como control. Los $F$ de primera etapa, casi todos inferiores a 10, muestran que al condicionar por shock y además partir por signo se agota la variación identificante: salvo el shock cambiario propio, los ERPT por signo no están identificados.
\end{tablenotes}
\end{threeparttable}
\end{table}